\documentclass[10pt,a4paper]{amsart}
\usepackage[margin=19mm]{geometry}
\usepackage{amsmath,amssymb,mathtools}
\usepackage[scr=rsfs,cal=euler]{mathalfa}
\usepackage{xfrac}
\usepackage[unicode,hidelinks,bookmarks=false]{hyperref}
\newcommand{\ie}{i.e.\ }
\newcommand{\cf}{cf.\ }
\newcommand{\resp}{respectively\ }
\newcommand{\Subsection}{\subsection}
\providecommand{\nobreakdash}{\nobreak\hbox{-}}
\setlength{\emergencystretch}{2em}
\multlinegap0pt
\usepackage{bm}
\usepackage{bbm}
\usepackage{dsfont}
\usepackage{xspace}
\usepackage{cancel}
\usepackage{mathtools}
\usepackage{amsthm}
\makeatletter
\@ifundefined{thrm}{\newtheorem{thrm}{Thrm}}{}
\@ifundefined{lemm}{\newtheorem{lemm}[thrm]{Lemma}}{}
\makeatother
\makeatletter
\expandafter\ifx\csname proof\endcsname\relax
\renewenvironment{proof}{\noindent{\bf Proof.}}{~~$\Box$}
\fi
\makeatother
\title[On the pointwise supremum of the set of copulas with a given]{On the pointwise supremum of the set of copulas with a given curvilinear section}
\author{Yao Ouyang, Yonghui Sun, Hua-Peng Zhang}
\date{Source: arXiv:2412.20629v4 (CC BY 4.0)}
\begin{document}
\maketitle

\noindent\emph{Source and licence.} On the pointwise supremum of the set of copulas with a given curvilinear section. Version arXiv:2412.20629v4 (CC BY 4.0), distributed under the Creative Commons Attribution 4.0 International licence, as recorded in the arXiv licence metadata for this exact version and shown on the abstract page https://arxiv.org/abs/2412.20629v4. Full rights notice: licenses/arxiv-2412.20629-CC-BY-4.0.txt.\par\smallskip

\noindent\textit{Editorial scope.} Counted unit: the Thrm whose source label is \texttt{None} in the article (source0.tex lines 352--354, source label \texttt{None}), delivered together with the article's own complete original proof of it (source0.tex lines 355--443, one \texttt{proof} environment of 89 lines). No mathematics is written, repaired or re-derived: statement and proof are the article's own text line for line. Required context retained uncounted: lemm-evofTV (lines 284--291). Every definition, display and label of those lines is retained unchanged. Omitted from this package and not redistributed: 1--283, 292--351, 444--991. Nothing inside a retained excerpt is omitted or altered: every retained body is delivered exactly as the article's lines are written, so no omission receipt of any kind accompanies this package. Citations: the retained excerpts carry no citation command at all, so no bibliography entry of the article is transcribed and no reference is redistributed. Reference adaptation: none was needed and none was made; every reference inside the delivered unit resolves inside it, so no unresolved reference or citation is produced. Printed slips kept literal: no printed word, symbol, hypothesis or number of the counted statement or of its proof was altered. Layout and class: the article's own class is \texttt{article} with packages including amsmath, amssymb, amsthm, latexsym, graphicx, xcolor, tikz, lineno; the wrapper is the lane's pinned \texttt{amsart} editorial wrapper at 10pt on A4, which restates the theorem-like environments the retained text needs and adds the article's own macro definitions that the retained text needs (none), with extra wrapper packages (bm, bbm, dsfont, xspace, cancel, mathtools). The delivered numbering is the unit's own. Assets: no retained span contains a figure environment, a graphics-inclusion command, a PDF-page inclusion, a table or a TikZ picture; no image file is copied into this package, the package declares no asset dependency and no image file is redistributed with it. Licence: version arXiv:2412.20629v4 of this article is distributed under the Creative Commons Attribution 4.0 International licence, as recorded in the arXiv licence metadata for this exact version and shown on the abstract page https://arxiv.org/abs/2412.20629v4. A full rights notice is delivered at licenses/arxiv-2412.20629-CC-BY-4.0.txt. Characters: every retained excerpt is pure ASCII, so the delivered engine, XeLaTeX with its default Latin Modern text font, reports no missing character.
\begin{lemm}\label{lemm-evofTV}
Let $\phi\in\Phi$ and $\Gamma_\phi\in\Delta_\phi$. Then the following inequalities hold for all $t_1, t_2\in[0, 1]$ with $t_1<t_2$:
\begin{itemize}\setlength{\itemindent}{10pt}
\item[{\rm (i)}] $\mathbb{V}_{t_1}^{t_2}(\hat{\Gamma}_\phi)\leq 2(\phi(t_2)-\phi(t_1))+ (\hat{\Gamma}_\phi(t_2)-\hat{\Gamma}_\phi(t_1))$;
\item[{\rm (ii)}] $\mathbb{V}_{t_1}^{t_2}(\tilde{\Gamma}_\phi)\leq 2(t_2-t_1)+(\tilde{\Gamma}_\phi(t_2)-\tilde{\Gamma}_\phi(t_1))$;
\item[{\rm (iii)}] $\mathbb{V}_{t_1}^{t_2}(\hat{\Gamma}_\phi)+ \mathbb{V}_{t_1}^{t_2}(\tilde{\Gamma}_\phi)\leq (t_2-t_1)+ (\phi(t_2)-\phi(t_1))$.
\end{itemize}
\end{lemm}

\begin{thrm}
Let $\phi\in\Phi$ and $\Gamma_\phi\in\Delta_\phi$. Then $C_1, C_2\in\mathbb{C}_{\Gamma_\phi}$.
\end{thrm}

\begin{proof}
We only prove that $C_1\in \mathbb{C}_{\Gamma_\phi}$ since the proof of $C_2\in \mathbb{C}_{\Gamma_\phi}$ is similar.
It is clear that $C_{1}(x, \phi(x))=\Gamma_\phi(x)$ for all $x\in[0, 1]$. It then remains to prove that $C_1$ is a copula.

\emph{Boundary condition:} For any $x\in[0, 1]$, we have
\[
\mathbb{V}^{x}_{\phi^{-1}(1)}(\hat{\Gamma}_{\phi})+\hat{\Gamma}_{\phi}(x)+\hat{\Gamma}_{\phi}(\phi^{-1}(1))
=\hat{\Gamma}_{\phi}(x)-\mathbb{V}^{1}_{x}(\hat{\Gamma}_{\phi})\leq 0.
\]
Hence, $f_1(x, 1)\geq x$, and thus $C_1(x, 1)=x$. It follows from Lemma~\ref{lemm-evofTV}(i) that
\[
\mathbb{V}^{1}_{\phi^{-1}(x)}(\hat{\Gamma}_{\phi})+\hat{\Gamma}_{\phi}(1)+\hat{\Gamma}_{\phi}(\phi^{-1}(x))\leq 2(\phi(1)-x)+2\hat{\Gamma}_{\phi}(1)=2(1-x).
\]
Hence, $f_1(1, x)\geq x$, and thus $C_1(1, x)=x$. Since $f_1$ is increasing, $C_1$ is also increasing, which implies that
$C_1(0, x)=C_1(x, 0)=0$.

\emph{$2$-increasingness:} Note that any rectangle $R=[a, b]\times [c, d]$ can be partitioned into at most three non-overlapping rectangles $R_i$ and $V(R)$ is the sum of all $V(R_i)$,
where each $R_i$ belongs to one of the three types of rectangles:

\noindent Type 1: $R\subset\{(x, y)\mid y\leq\phi(x)\}$, \emph{i.e.}, $R$ is below the curve $\{(x, \phi(x))\mid x\in [0, 1]\}$;

\noindent Type 2: $R\subset\{(x, y)\mid y\geq\phi(x)\}$, \emph{i.e.}, $R$ is above the curve $\{(x, \phi(x))\mid x\in [0, 1]\}$;

\noindent Type 3: $R=[x_1, x_2]\times [\phi(x_1), \phi(x_2)]$, \emph{i.e.}, $R$ has two vertices on the curve $\{(x, \phi(x))\mid x\in [0, 1]\}$.

\noindent We only need to prove that $V(R)\geq 0$ for each of the three types of rectangles $R=[x_1, x_2]\times [y_1, y_2]$.

\noindent Type 1: $R=[x_1, x_2]\times[y_1, y_2]\subset\{(x, y)\mid y\leq\phi(x)\}$.

In this case, we have $C_{1}(x, y)=\min\{y, f_1(x, y)\}$ for any $(x, y)\in R$.

If $C_1(x_1, y_1)=y_1$, then it follows from the increasingness of $C_1$ that $V(R)=C_1(x_2, y_2)-C_1(x_1, y_2)\geq 0$.

For the case that $C_1(x_1, y_1)=f_1(x_1, y_1)$, there are two subcases to be considered:
\begin{itemize}%\setlength{\itemindent}{5pt}
\item[\rm (a)] $C_1(x_2, y_2)=y_2$.

In this subcase, we have
\[
V(R)\geq f_1(x_1, y_1)+y_2-y_1-f_1(x_1, y_2)\geq 0,
\]
where the last inequality is due to the fact that $f_1$ is $1$-Lipschitz.

\item[\rm (b)] $C_1(x_2, y_2)=f_1(x_2, y_2)$.

In this subcase, we have
\[
V(R)\geq f_1(x_1, y_1)+f_1(x_2, y_2)-f_1(x_1, y_2)-f_1(x_2, y_1)=0.
\]
\end{itemize}

\noindent Type 2: $R=[x_1, x_2]\times[y_1, y_2]\subset\{(x, y)\mid y\geq\phi(x)\}$.

If $C_1(x_1, y_1)=x_1$ or $C_1(x_1, y_1)=y_1$, then the increasingness of $C_1$ implies that $V(R)\geq 0$. So we suppose that
$C_1(x_1, y_1)=f_1(x_1, y_1)$. There are three subcases to be considered:
\begin{itemize}%\setlength{\itemindent}{5pt}
\item[\rm (a)] $C_1(x_2, y_2)=x_2$.

In this subcase, we have
\[
V(R)\geq f_1(x_1, y_1)+x_2-x_1-f_1(x_2, y_1)\geq 0.
\]

\item[\rm (b)] $C_1(x_2, y_2)=y_2$.

In this subcase, we have
\[
V(R)\geq f_1(x_1, y_1)+y_2-y_1-f_1(x_1, y_2)\geq 0.
\]

\item[\rm (c)] $C_1(x_2, y_2)=f_1(x_2, y_2)$.

In this subcase, we have
\[
V(R)\geq f_1(x_1, y_1)+f_1(x_2, y_2)-f_1(x_1, y_2)-f_1(x_2, y_1)=0.
\]
\end{itemize}

\noindent Type 3: $R=[x_1, x_2]\times [\phi(x_1), \phi(x_2)]$.

In this case, we have $C_1(x_1, \phi(x_1))=\Gamma_\phi(x_1)$ and $C_1(x_2, \phi(x_2))=\Gamma_\phi(x_2)$. Hence,
\begin{eqnarray*}
V(R)&\geq& \Gamma_\phi(x_1)+\Gamma_\phi(x_2)-f_1(x_1, \phi(x_2))-f_1(x_2, \phi(x_1))\\
&=& \Gamma_\phi(x_1)+\Gamma_\phi(x_2)-x_1+\frac{1}{2}\Big(\mathbb{V}^{x_1}_{x_2}(\hat{\Gamma}_{\phi})+\hat{\Gamma}_{\phi}(x_1)+\hat{\Gamma}_{\phi}(x_2)\Big)\\
&&- x_2+ \frac{1}{2}\Big(\mathbb{V}^{x_2}_{x_1}(\hat{\Gamma}_{\phi})+\hat{\Gamma}_{\phi}(x_1)+\hat{\Gamma}_{\phi}(x_2)\Big)\\
&=& \Gamma_\phi(x_1)+\Gamma_\phi(x_2)-\Gamma_\phi(x_1)-\Gamma_\phi(x_2)=0.
\end{eqnarray*}
We conclude that $C_1\in \mathbb{C}_{\Gamma_\phi}$.
\end{proof}

\end{document}
